% =============================================================================
%  Sección 1: Introducción
% =============================================================================

\section{Introducción}
\label{sec:introduccion}

El urbanismo en México no puede analizarse sin un estudio profundo de su historia y evolución a lo largo del siglo XX
y XXI. Desde la Revolución Mexicana acontecida en 1910 y que culminó oficialmente con la promulgación de la Constitución
Política de los Estados Unidos Mexicanos en 1917. Sin embargo, México como país no despegaría hasta
la época conocida como el \termino{Cardenismo}, el cual tuvo como figura principal al general Lázaro Cárdenas del Río, quien
influyó en la consolidación del \termino{Partido de la Revolución Mexicana (PRM)}, anteriormente \termino{Partido Nacional Revolucionario (PNR)}
y posteriormente conocido como \termino{Partido Revolucionario Institucional (PRI)} el cual sería el arquitecto del sistema económico 
que sustentaría al país durante dos épocas clave \termino{El Milagro Mexicano} y la consolidación del Neoliberalismo como corriente
económica en México. A esta etapa se le puede llamar \termino{La dictadura Perfecta} \cite{vargasllosa_dictadura}. Sin embargo, la sociedad y su manera 
de concebir su vida personal, profesional y académica atravesaría por diferentes cambios en su cultura y forma de concebir la realidad 
del México Moderno y Contemporáneo. En este ensayo se explorarán 5 puntos particulares, pero que se resumen fácilmente a través de 3 ejes
principales: \textbf{Sociedad, Gobierno y Transporte}.


\subsection{Contexto metropolitano en el siglo XXI}
\label{ssec:contexto-metropolitano}

Analizar la \zmvm, puede convertirse en una tarea titánica y con ánimos de nunca terminar. Para ello se puede usar la técnica conocida como: 
\termino{Divide y Vencerás}. Entender la urbe mexicana a través de sus cambios generacionales nos proporciona una visión más exacta de la 
sinergia y convivencia que tienen sus habitantes dependiendo de su rango de edad; \termino{Generación Silenciosa (1928-1945)}, 
\termino{Baby Boomer (1946-1964)}, \termino{Generación X (1965-1980)}, \termino{Milenials (1981-1996)}, \termino{Generación Z (1997-2012)} y 
\termino{Generación Alfa (2010-2020)} \cite{pew_generations_2019}. Esta división no es casualidad, pues cada una atravesó diferentes cambios sociales 
y económicos que las marcaron e hicieron únicas a cada una. Sin embargo, no se puede negar que el país atraviesa por un periodo de \termino{Coyuntura Social} 
y \termino{Coyuntura Política} ocasionada por el crecimiento demográfico desbordado y la planeación insuficiente de sus respectivas administraciones. 
A la fecha de este ensayo (Mayo 2026), las generaciones Milenials y Alfa son las que cargan en sus hombros la economía y fuerza bruta del país por su edad 
productiva, ya que se encuentran dentro del rango conocido como \termino{\textbf{Población Económicamente Activa (PEA)}}, el cual va de los 18 a los 65 años de edad. 
A pesar de ser de las generaciones más estudiadas de la historia \cite{pew_generations_2019}, sus ingresos y posibilidades de adquirir bienes inmobiliarios o 
créditos automotrices son infinitamente menores a los de sus predecesores. Si la economía de la \termino{PEA} no puede sostener una calidad de vida aceptable 
según métricas y parámetros de la ONU \cite{undp_hdr_2024}, ocurren los periodos de crisis que atañen el problema metropolitano de este ensayo. 
Entender que esto no es un problema aislado y que no existe un solo culpable (ni generacional, ni político, ni económico) nos da un mejor panorama de cómo
diagnosticarlo, aunque la cura aún esté en proceso de creación. 


\subsection{Objetivo del ensayo}
\label{ssec:objetivo-ensayo}

El presente ensayo tiene por objetivo dar un análisis cualitativo en su mayoría, presentando datos cuantitativos de la investigación del tema de tesis:
\textit{Transportes Anillares y su Importancia en la Ciudad de México y el área metropolitana} \cite{galigaribaldi_tesis} a través de las herramientas de
análisis \textit{Apimetro \cite{galigaribaldi_apimetro}} y \textit{Modelo de la Higuera Evanescente o VFTModel \cite{galigaribaldi_vftmodel}}. Centrado en 5 puntos:

\begin{itemize}
    \item Analizar la disputa y correlación que tienen las Ciencias Sociales como la sociología y la filosofía con las Ciencias Exactas como las Matemáticas y la Computación.
    \item Analizar y dar una visión cualitativa del Estado como actor económico a través de sus modelos de Gobierno.
    \item Analizar y dar una visión cualitativa de la sociedad mexicana como resultado de la gestión pública dentro de la \zmvm.
    \item Analizar cómo la ciudad y las necesidades económicas han orillado a sus habitantes a un \termino{Estándar} de personalidad.
    \item Conectar los puntos sociales, económicos y de gobierno en un espacio terciario como lo es el Transporte Público.
\end{itemize}
